\NormalFont\ShortTitle{Hebreus}
{\MT Carta aos Hebreus

\par }\ChapOne{1}{\PP \VerseOne{1}Deus falou antigamente muitas vezes, e de muitas maneiras, aos ancestrais
\FTNT{*}{{\FR 1:1 }
lit. pais
} pelos profetas.
\VS{2}Mas nestes últimos falou a nós por meio do
{\ADD{seu}} Filho, a quem constituiu por herdeiro de todas as coisas, e pelo qual também fez o universo.
\FTNT{†}{{\FR 1:2 }
o universo
lit. as eras – o termo grego literalmente indicava longos períodos de tempo, mas também podia indicar lugares, isto é, o mundo ou mundos
}
\VS{3}{\ADD{O Filho}} é o resplendor da
{\ADD{sua}} glória, a expressa imagem da sua pessoa, e sustenta todas as coisas pela sua palavra de poder.
\FTNT{‡}{{\FR 1:3 }
sua palavra de poder
lit. palavra de seu poder
} E, depois de fazer por si mesmo a purificação dos nossos pecados, sentou-se à direita da Majestade nas alturas;
\VS{4}e se tornou muito superior aos anjos, assim como herdou um nome mais excelente que eles.
\VS{5}Pois a qual dos anjos
{\ADD{Deus}} jamais disse: “Tu és meu Filho, eu hoje te gerei”,
{\RQ{Salmos 2:7
}} E em outra vez: “Eu lhe serei por Pai, e ele me será por Filho”?
{\RQ{2 Samuel 7:14, 1 Crônicas 17:13
}}
\VS{6}E outra vez, quando introduz o primogênito ao mundo, diz: E todos os anjos de Deus o adorem.
{\RQ{Salmos 97:7
}}
\VS{7}Quanto aos anjos, porém, ele diz: Ele faz de seus anjos espíritos; e de seus trabalhadores, labareda de fogo;
{\RQ{Salmos 104:4
}}
\VS{8}Mas ao Filho
{\ADD{diz}} : Ó Deus, o teu trono é para todo o sempre. Cetro de equidade é o cetro do teu Reino.
\VS{9}Amaste a justiça, e odiaste a transgressão; por isso, Deus, o teu Deus, te ungiu com óleo de alegria mais que os teus companheiros.
{\RQ{Salmos 45:6-7
}}
\VS{10}E: Tu Senhor, no princípio fundaste a terra, e os céus são obras de tuas mãos.
\VS{11}Eles perecerão, mas tu continuas. Todos eles como roupa se envelhecerão;
\VS{12}como uma manta tu os envolverás, e serão mudados; porém tu és o mesmo, e os teus anos não cessarão.
{\RQ{Salmos 102:25-27
}}
\VS{13}E a qual dos anjos ele jamais disse: Senta-te à minha direita, até que eu ponha os teus inimigos como estrado de teus pés?
{\RQ{Salmos 110:1
}}
\VS{14}Por acaso não são todos eles espíritos servidores, enviados para auxílio dos que herdarão a salvação?

\par }\Chap{2}{\PP \VerseOne{1}Portanto, devemos prestar atenção com muito mais empenho para as coisas que já ouvimos, para que nunca venhamos a cometer deslizes.
\VS{2}Pois, se a palavra pronunciada pelos anjos foi confirmada, e toda transgressão e desobediência recebeu justa retribuição,
\VS{3}como nós escaparemos, se descuidarmos de tão grande salvação? Ela, que começou a ser anunciada pelo Senhor, foi-nos confirmada pelos que ouviram.
\VS{4}Deus também deu testemunho com eles, por meio de sinais, milagres, várias maravilhas, e distribuições do Espírito Santo, segundo a sua vontade.
\VS{5}Pois não foi aos anjos que ele sujeitou o mundo futuro, do qual estamos falando.
\VS{6}Mas em certo lugar alguém testemunhou, dizendo: O que é o homem, para que te lembres dele? Ou o filho do homem, para que o visites?
\VS{7}Tu o fizeste um pouco menor que os anjos, de glória e de honra o coroaste, e o estabeleceste sobre as obras das tuas mãos;
\VS{8}Sujeitaste todas as coisas debaixo de seus pés.Pois ao lhe sujeitar todas as coisas, nada deixou que não fosse sujeito a ele; porém agora ainda não vemos todas as coisas sendo sujeitas a ele.
{\RQ{Salmos 8:4-6
}}
\VS{9}Porém vemos coroado de glória e de honra, por causa do sofrimento de morte, aquele Jesus, que havia sido feito um pouco
\FTNT{*}{{\FR 2:9 }
um pouco
ou: por um pouco [de tempo]
} menor que os anjos, a fim de que, pela graça de Deus, experimentasse a morte por todos.
\VS{10}Pois convinha a aquele, para quem e por quem são todas as coisas, que, quando trouxesse muitos filhos à glória, aperfeiçoasse por meio de aflições o Autor
\FTNT{†}{{\FR 2:10 }
Autor
ou: Príncipe, Primeiro
} da salvação deles.
\VS{11}Porque tanto o que santifica como os que são santificados, todos são provenientes de um; por isso ele não se envergonha de lhes chamar de irmãos,
\VS{12}dizendo: Anunciarei o teu nome aos meus irmãos, cantarei louvores a ti no meio da congregação.
{\RQ{Salmos 22:22
}}
\VS{13}E outra vez: Nele confiarei.
{\RQ{Isaías 8:17
}} E outra vez: Eis-me aqui, com os filhos que Deus me deu.
{\RQ{Ref. Isaías 8:18
}}
\VS{14}Portanto, uma vez que os filhos compartilham da carne e do sangue, também ele, semelhantemente, participou das mesmas coisas, a fim de pela morte aniquilar ao que tinha o poder da morte, isto é, o diabo;
\VS{15}e libertar a todos os que, por medo da morte, estavam sujeitos à servidão durante a vida toda.
\VS{16}Pois, evidentemente, não é aos anjos que ele ajuda, mas sim, ajuda à descendência
\FTNT{‡}{{\FR 2:16 }
descendência
lit. semente
} de Abraão.
\VS{17}Por isso era necessário que em tudo ele se tornasse semelhante aos irmãos, para ser misericordioso e fiel Sumo Sacerdote nas coisas
{\ADD{relativas}} a Deus, para fazer sacrifício de perdão
\FTNT{§}{{\FR 2:17 }
sacrifício de perdão
equiv. reconciliação, propiciação. A palavra grega se refere ao apaziguamento da fúria divina através de um sacrifício
} dos pecados do povo.
\VS{18}Pois, naquilo em que ele mesmo sofreu ao ser tentado, ele pode socorrer aos que são tentados.

\par }\Chap{3}{\PP \VerseOne{1}Portanto, santos irmãos, participantes do chamado celestial, considerai o Apóstolo e Sumo Sacerdote da nossa declaração de fé, Cristo Jesus.
\VS{2}Ele é fiel àquele que o constituiu, assim como Moisés também foi em toda a casa de Deus.
\FTNT{*}{{\FR 3:2 }
casa de Deus
lit. sua casa. Também no v. 5
}
\VS{3}Pois ele é estimado como digno de maior glória do que Moisés, assim como mais honra do que a casa tem quem a construiu.
\VS{4}Porque toda casa é construída por alguém; mas aquele que construiu todas as coisas é Deus.
\VS{5}De fato, Moisés foi fiel como servo em toda a casa de Deus, para testemunho das coisas que haviam de ser ditas;
\VS{6}mas Cristo
{\ADD{é fiel}} como Filho sobre a sua casa. E nós somos sua casa, se até o fim mantivermos firmes a confiança e o orgulho da
{\ADD{nossa}} esperança.
\FTNT{†}{{\FR 3:6 }
orgulho da [nossa] esperança
i.e. a esperança que nos orgulhamos de ter
}
\VS{7}Portanto, como diz o Espírito Santo: Hoje, se ouvirdes a sua voz,
\VS{8}não endureçais os vossos corações, como na provocação,
\FTNT{‡}{{\FR 3:8 }
provocação
ou: rebelião – também no v. 15
} no dia da tentação no deserto;
\VS{9}onde os vossos pais
\FTNT{§}{{\FR 3:9 }
pais
i.e. ancestrais
} me tentaram, me provaram, e viram as minhas obras por quarenta anos.
\VS{10}Por isso me indignei contra essa geração, e disse: “Eles sempre se desviam nos corações, e não conheceram os meus caminhos”;
\VS{11}Então jurei na minha ira: “Eles não entrarão no meu repouso”.
{\RQ{Salmos 95:7-11
}}
\VS{12}Cuidado, irmãos, para que nunca haja em algum de vós coração mau e infiel,
\FTNT{*}{{\FR 3:12 }
infiel
ou: incrédulo. Lit. coração mau de infidelidade (ou incredulidade). No grego, a fé fidelidade eram representados pela mesma palavra, ocorrendo o mesmo para os opostos
} que se afaste do Deus vivo.
\VS{13}Em vez disso, encorajai-vos uns aos outros a cada dia, durante o tempo chamado “hoje”; a fim de que nenhum de vós se endureça pelo engano do pecado.
\VS{14}Pois nós temos nos tornado participantes de Cristo, se, de fato, mantivermos firme a confiança do princípio
\FTNT{†}{{\FR 3:14 }
confiança do princípio
lit. princípio da confiança
} até o fim;
\VS{15}como é dito: Hoje, se ouvirdes a sua voz, não endureçais os vossos corações, como na provocação.
\VS{16}Porque alguns, mesmo ouvindo, provocaram-no
{\ADD{à ira}} ,
\FTNT{‡}{{\FR 3:16 }
provocaram-no [à ira]
ou: rebelaram-se
} mas não todos aqueles que saíram do Egito por meio de Moisés.
\FTNT{§}{{\FR 3:16 }
Porque alguns … Moisés
Muitas traduções interpretam a frase na forma de perguntas: “Pois quais foram os que, mesmo ouvindo, provocaram-no [à ira] ? Não foram todos aqueles que saíram do Egito por meio de Moisés?
}
\VS{17}E contra quais ele se indignou por quarenta anos? Não foi contra os que pecaram, cujos corpos caíram no deserto?
\VS{18}E a quem ele jurou que não entrariam no seu repouso, senão aos que foram desobedientes?
\VS{19}Assim vemos que não puderam entrar por causa da incredulidade.

\par }\Chap{4}{\PP \VerseOne{1}Temamos, pois, para que, havendo sido deixada a promessa de entrar no seu repouso, não pareça que algum de vós tenha ficado para trás.
\VS{2}Pois também a nós foi pregado o evangelho, assim como a eles; mas a palavra da pregação nada lhes aproveitou, porque não estava combinada com a fé naqueles que
{\ADD{a}} ouviram.
\VS{3}Porque nós, os que cremos, entramos no repouso, como ele disse: Então jurei na minha ira: “Eles não entrarão no meu repouso”; ainda que as
{\ADD{suas}} obras já estivessem terminadas desde a fundação do mundo.
\VS{4}Pois, em certo lugar, ele disse assim a respeito do
{\ADD{dia}} sétimo: E Deus repousou no sétimo dia de todas as suas obras.
{\RQ{Gênesis 2:2
}}
\VS{5}E, outra vez neste
{\ADD{texto}} : Não entrarão no meu repouso.
\VS{6}Portanto, uma vez que resta alguns entrarem no repouso,
\FTNT{*}{{\FR 4:6 }
no repouso
lit. nele
} e que aqueles a quem primeiramente foi pregado o evangelho não entraram por causa da desobediência,
\VS{7}outra vez ele determina um certo dia, o “hoje”, dizendo por meio de Davi muito tempo depois (como já foi dito): Hoje, se ouvirdes a sua voz, não endureçais os vossos corações.
\VS{8}Pois se Josué tivesse lhes dado repouso, ele não teria falado depois a respeito de outro dia.
\VS{9}Portanto, ainda resta um repouso como o do sábado para o povo de Deus.
\VS{10}Pois aquele que entrou no seu repouso também repousou das suas obras, assim como Deus das suas.
\VS{11}Procuremos, pois, entrar naquele repouso; para que ninguém caia no mesmo exemplo de desobediência.
\VS{12}Porque a palavra de Deus é viva e eficaz, e mais afiada do que toda espada de dois gumes, e que penetra até a divisão da alma e do espírito, das juntas e das medulas, e é apta para julgar os pensamentos e intenções do coração.
\VS{13}E não há criatura alguma encoberta diante dele; pelo contrário, todas as coisas estão nuas e expostas aos olhos daquele a quem temos de prestar contas.
\VS{14}Visto que temos um grande Sumo Sacerdote, Jesus, o Filho de Deus, que penetrou nos céus, mantenhamos firme a
{\ADD{nossa}} declaração de fé.
\VS{15}Pois não temos um Sumo Sacerdote que não possa se compadecer das nossas fraquezas; mas, sim, um que foi provado em tudo, conforme a
{\ADD{nossa}} semelhança,
{\ADD{porém}} sem pecado.
\VS{16}Portanto, acheguemo-nos com confiança ao trono da graça, para que possamos receber misericórdia e encontrar graça para sermos socorridos
\FTNT{†}{{\FR 4:16 }
para sermos socorridos
lit. para socorro
} no tempo adequado.

\par }\Chap{5}{\PP \VerseOne{1}Pois todo sumo sacerdote tomado dentre os homens é constituído em prol das pessoas nas coisas
{\ADD{relativas}} a Deus, para apresentar tanto ofertas como sacrifícios pelos pecados;
\VS{2}e ele é capaz de se compadecer dos ignorantes e dos errados, pois ele mesmo também está rodeado de fraqueza.
\VS{3}E, por causa desta
{\ADD{fraqueza}} ,ele tem o dever de apresentar oferenda, tanto pelos pecados do povo, como também de si mesmo.
\VS{4}E ninguém toma para si esta honra, a não ser aquele que é chamado por Deus, como Arão.
\VS{5}Assim também Cristo não glorificou a si mesmo para se fazer Sumo Sacerdote; mas sim, àquele que lhe disse: Tu és meu Filho, eu hoje te gerei.
\VS{6}Como também diz em outro
{\ADD{texto}} : Tu és Sacerdote para sempre, segundo a ordem de Melquisedeque.
{\RQ{Salmos 110:4
}}
\VS{7}Ele, nos dias de seu corpo carnal, ofereceu com grande clamor e lágrimas, tanto orações como súplicas para aquele que podia o livrar da morte, e foi ouvido por causa da devoção.
\FTNT{*}{{\FR 5:7 }
devoção
ou: temor
}
\VS{8}Ainda que ele era o Filho, aprendeu a obediência por meio das coisas que sofreu.
\VS{9}E, depois de ter sido consumado, ele foi feito autor da eterna salvação a todos os que lhe obedecem;
\VS{10}e nomeado por Deus como Sumo Sacerdote, segundo a ordem de Melquisedeque.
\VS{11}Nosso comentário a respeito disso é extenso e difícil de explicar, porque vos tornastes negligentes para ouvir.
\VS{12}Pois, pelo tempo, vós já devíeis ser mestres; porém novamente tendes necessidade de serdes ensinados os princípios básicos das palavras de Deus; e vos fizestes como se precisásseis de leite, e não de alimento sólido.
\VS{13}Pois qualquer um que depende
\FTNT{†}{{\FR 5:13 }
depende
lit. participa
} de leite é inexperiente na palavra da justiça; porque
{\ADD{ainda}} é um bebê.
\VS{14}Mas o alimento sólido é para os maduros, os quais, pelo costume, têm os sentidos treinados para discernirem tanto o bem como o mal.

\par }\Chap{6}{\PP \VerseOne{1}Portanto deixemos os tópicos elementares da doutrina de Cristo, e prossigamos adiante até a maturidade, sem lançarmos de novo o fundamento da conversão de obras mortas, e da fé em Deus,
\VS{2}da doutrina dos batismos, da imposição das mãos, da ressurreição dos mortos, e do juízo eterno.
\VS{3}E isso faremos, se Deus permitir.
\VS{4}Pois é impossível que os que uma vez foram iluminados, e experimentaram o dom celestial, e foram feitos participantes do Espírito Santo,
\VS{5}e experimentaram a boa palavra de Deus, e os poderes do mundo futuro;
\VS{6}e então caíram, outra vez renová-los para o arrependimento, porque estão novamente crucificando o Filho de Deus para si mesmos, e o expondo à vergonha pública.
\FTNT{*}{{\FR 6:6 }
conversão
ou: arrependimento
}
\VS{7}Pois a terra que absorve a chuva que com frequência vem sobre ela, e produz vegetação proveitosa para aqueles por quem e lavrada, recebe a bênção de Deus.
\VS{8}Contudo, se ela produz espinhos e cardos, é reprovada, e perto está da maldição. O seu fim é ser queimada.
\VS{9}Mas, ainda que estejamos falando assim, quanto a vós mesmos, amados, confiamos melhores coisas, relacionadas à salvação.
\VS{10}Pois Deus não é injusto para se esquecer da vossa obra, e do trabalho do amor que mostrastes para com o nome dele, quando servistes aos santos, e continuais a servir.
\VS{11}Mas desejamos que cada um de vós mostre o mesmo empenho até o fim, para que a
{\ADD{vossa}} esperança seja completa;
\FTNT{†}{{\FR 6:11 }
para que a [vossa] esperança seja completa
lit. para a completude (ou: completa certeza) da esperança
}
\VS{12}a fim de que não vos torneis negligentes, mas sim, imitadores dos que pela fé e pela paciência herdam as promessas.
\VS{13}Pois Deus, quando prometeu a Abraão, como ninguém maior havia por quem jurar, jurou por si mesmo,
\VS{14}dizendo: Certamente muito te abençoarei,
\FTNT{‡}{{\FR 6:14 }
muito te abençoarei
lit. abençoando te abençoarei
} e muito te multiplicarei.
\FTNT{§}{{\FR 6:14 }
muito te multiplicarei
lit. multiplicando te multiplicarei
}
{\RQ{Gênesis 22:17
}}
\VS{15}E assim, esperando pacientemente,
{\ADD{Abraão}} obteve a promessa.
\VS{16}Pois as pessoas juram por alguém maior, e o juramento como confirmação é o fim de todo conflito entre elas.
\VS{17}Assim Deus, querendo mostrar mais abundantemente a imutabilidade de sua intenção aos herdeiros da promessa, garantiu com um juramento;
\VS{18}a fim de que, por meio de duas coisas imutáveis, nas quais é impossível que Deus minta, nós, que buscamos refúgio, tenhamos encorajamento para agarrar a esperança posta diante
{\ADD{de nós}} .
\VS{19}Temos essa esperança
\FTNT{*}{{\FR 6:19 }
essa esperança
incluída explicitamente por questão de fluência. No texto grego referida através de pronome relativo
} como uma segura e firme âncora da alma, que entra ao interior, além do véu,
\VS{20}onde,
{\ADD{como}} precursor, Jesus entrou para nosso benefício, pois ele se tornou Sumo Sacerdote para sempre, segundo a ordem de Melquisedeque.

\par }\Chap{7}{\PP \VerseOne{1}Porque este Melquisedeque, rei de Salém, sacerdote do Deus Altíssimo, o qual se encontrou com Abraão, que estava retornando da matança dos reis, e o abençoou;
\VS{2}a quem também Abraão deu o dízimo de tudo; primeiramente significa Rei de justiça, e depois também rei de Salém, que é, rei de paz;
\VS{3}sem pai, sem mãe, sem genealogia, sem ter princípio de dias, nem fim de vida; porém, sendo feito semelhante ao Filho de Deus, permanece sacerdote para sempre.
\VS{4}Considerai, pois, como ele era grande, a quem até o patriarca Abraão deu o dízimo dos despojos.
\VS{5}E, realmente, os que dentre os filhos de Levi recebem o sacerdócio têm ordem, segundo a Lei, de receberem os dízimos do povo, isto é, dos seus irmãos, mesmo sendo eles também descendentes de Abraão.
\FTNT{*}{{\FR 7:5 }
sendo eles descendentes de Abraão
lit. tendo eles [também] saído dos lombos de Abraão.
}
\VS{6}Mas aquele que não é contado na genealogia deles recebeu dízimos de Abraão, e abençoou ao que tinha as promessas.
\VS{7}Ora, sem contradição alguma, o menor é abençoado pelo maior.
\VS{8}Em um caso, homens mortais recebem dízimos; mas no outro, aquele de quem se dá testemunho de que vive.
\VS{9}E, por assim dizer, até Levi, que recebe os dízimos, pagou dízimos através de Abraão;
\VS{10}pois ele ainda estava no corpo
\FTNT{†}{{\FR 7:10 }
corpo
lit. lombos
} de
{\ADD{seu}} ancestral
\FTNT{‡}{{\FR 7:10 }
ancestral
lit. pai
} quando Melquisedeque se encontrou com ele.
\VS{11}Portanto, se a perfeição tivesse sido de fato pelo sacerdócio Levítico (pois sob ele o povo recebeu a Lei), que mais necessidade havia de se levantar outro Sacerdote segundo a ordem de Melquisedeque, e não ser chamado segundo a ordem de Arão?
\VS{12}Pois, ao se mudar o sacerdócio, necessariamente também se faz mudança de Lei.
\VS{13}Porque aquele de quem estas coisas são ditas pertence a outra tribo, da qual ninguém serviu ao altar;
\VS{14}visto ser evidente que o nosso Senhor é procedente de Judá, tribo da qual Moisés nada falou a respeito de sacerdócio.
\VS{15}E
{\ADD{isso}} ainda é muito mais evidente se, à semelhança de Melquisedeque, levanta-se outro sacerdote,
\VS{16}que foi constituído, não conforme a Lei de um mandamento carnal, mas sim, conforme o poder de uma vida indestrutível.
\VS{17}Pois
{\ADD{assim}} ele dá testemunho: Tu és Sacerdote para sempre, segundo a ordem de Melquisedeque.
\VS{18}Porque há uma revogação do mandamento anterior, por causa da sua fraqueza e inutilidade,
\VS{19}pois a Lei não tornou perfeito nada. Porém, uma esperança melhor é introduzida, e por meio dela nos aproximamos de Deus.
\VS{20}E isso não foi
{\ADD{feito}} sem juramento (pois os outros se tornaram sacerdotes sem juramento,
\VS{21}mas ele
{\ADD{foi estabelecido}} com um juramento daquele que lhe disse: O Senhor jurou, e não se arrependerá: Tu és Sacerdote para sempre, segundo a ordem de Melquisedeque).
\VS{22}Assim, Jesus foi feito fiador de um pacto
\FTNT{§}{{\FR 7:22 }
pacto
ou: testamento
} ainda melhor.
\VS{23}Dos outros, são muitos os que se tornaram sacerdotes, pois pela morte foram impedidos de continuar;
\VS{24}mas ele, porque permanece para sempre, tem um sacerdócio definitivo.
\VS{25}Portanto, ele também pode salvar de maneira completa os que se aproximam de Deus por meio dele, visto que ele vive para sempre para interceder por eles.
\VS{26}Pois nos era conveniente tal Sumo Sacerdote: santo, inocente, incontaminado, separado dos pecadores, e feito mais elevado que os céus;
\VS{27}que não precisasse, como os sumos sacerdotes, de oferecer sacrifícios diariamente, primeiramente pelos seus próprios pecados, e depois pelos do povo. Pois ele fez isto de uma vez por todas quando ofereceu a si mesmo.
\VS{28}Porque a Lei constitui por sumos sacerdotes homens que têm fraqueza; mas a palavra do juramento, que
{\ADD{veio}} depois da Lei,
{\ADD{constitui}} o Filho, que se tornou perfeito para sempre.

\par }\Chap{8}{\PP \VerseOne{1}Ora, o ponto-chave do que falamos é que temos um Sumo Sacerdote que está sentado à direita do trono da Majestade nos céus,
\VS{2}servidor
\FTNT{*}{{\FR 8:2 }
servidor
ou: ministro
} do Santuário, e do verdadeiro Tabernáculo, que Senhor ergueu, e não o homem.
\VS{3}Pois todo sumo sacerdote é constituído para apresentar tanto ofertas como sacrifícios; por isso era necessário que também este tivesse alguma coisa a oferecer.
\VS{4}Mas, se ele
{\ADD{ainda}} estivesse na terra, nem sequer seria sacerdote, já que
{\ADD{ainda}} há os sacerdotes que apresentam as ofertas segundo a Lei.
\VS{5}Esses servem a um esboço, uma sombra das coisas celestiais, como Moisés foi divinamente avisado, antes de construir o Tabernáculo. Porque
{\ADD{Deus}} disse: Olha, faze tudo conforme o modelo que te foi mostrado no monte.
{\RQ{Êxodo 25:40
}}
\VS{6}Mas agora
{\ADD{Jesus}} obteve um ofício
\FTNT{†}{{\FR 8:6 }
ofício
ou: ministério
} mais relevante,
\FTNT{‡}{{\FR 8:6 }
relevante
ou: excelente
} como é também Mediador de um melhor pacto, que está firmado
\FTNT{§}{{\FR 8:6 }
firmado
ou: legalmente estabelecido
} em melhores promessas.
\VS{7}Pois, se aquele primeiro tivesse sido infalível, nunca haveria se buscado lugar para o segundo.
\VS{8}Porque, enquanto
{\ADD{Deus}} repreendia
\FTNT{*}{{\FR 8:8 }
repreendia os do povo
lit. repreendia-os
} os do povo,
\FTNT{†}{{\FR 8:8 }
povo
lit. casa. O termo tradicionalmente traduzido como “casa” pode incluir vários significados, entre eles, o de povo. Neste caso, os povos formados pelos descendentes de Israel e de Judá
} disse-lhes: Eis que vêm dias, diz o Senhor, que, sobre o povo de Israel e sobre o povo de Judá, estabelecerei um novo pacto;
\VS{9}não conforme o pacto que fiz com os seus ancestrais
\FTNT{‡}{{\FR 8:9 }
ancestrais
lit. pais
} no dia em que eu os tomei pelas mãos, para os tirar da terra do Egito. Pois não permaneceram naquele meu pacto, e parei de lhes dar atenção, diz o Senhor.
\VS{10}Este, pois, é o pacto que farei com o povo de Israel depois daqueles dias, diz o Senhor: Porei minhas leis nas mentes deles, e as escreverei em seus corações. Eu serei o Deus deles, e eles serão o meu povo.
\VS{11}E não terão de ensinar cada um ao seu próximo, ou ao seu irmão, dizendo: “Conhece ao Senhor”; porque todos me conhecerão, desde o menor deles até o maior.
\VS{12}Pois serei misericordioso com as suas injustiças, e não mais me lembrarei dos seus pecados e das suas transgressões.
{\RQ{Jeremias 31:31-34
}}
\VS{13}Quando ele disse: “Novo
{\ADD{pacto}} ”, ele tornou o primeiro ultrapassado. Ora, o que se torna ultrapassado e envelhece está perto de desaparecer.

\par }\Chap{9}{\PP \VerseOne{1}O primeiro
{\ADD{pacto}} também tinha ordenanças de culto, e o santuário terrestre.
\VS{2}Pois um tabernáculo foi preparado, o primeiro, em que havia o candelabro, a mesa, e os pães da proposição. Esse é chamado o Santo Lugar.
\VS{3}Mas após o segundo véu estava o tabernáculo que se chama Santo dos Santos;
\FTNT{*}{{\FR 9:3 }
Santo dos Santos
ou: Santíssimo Lugar
}
\VS{4}que tinha o incensário de ouro, e a arca do pacto, toda coberta de ouro. Nela estavam: o vaso de ouro contendo o maná, a vara de Arão que florescera, e as tábuas do pacto.
\VS{5}E acima dela, os querubins de glória, que faziam sombra ao propiciatório. Acerca dessas coisas não
{\ADD{é oportuno}} agora falar em detalhes.
\VS{6}Ora, estando estas coisas assim preparadas, os sacerdotes entram a todo tempo no primeiro tabernáculo, para cumprirem as atividades de culto.
\VS{7}Mas no segundo, somente o sumo sacerdote, uma vez por ano, não sem sangue, que oferece por si mesmo, e pelos pecados de ignorância do povo.
\VS{8}Desta maneira, o Espírito Santo dá a entender que o caminho para o Santuário
\FTNT{†}{{\FR 9:8 }
Há controvérsia se a referência é ao lugar Santo, ou ao Santo dos Santos
} não havia sido revelado enquanto o primeiro Tabernáculo ainda estava de pé.
\VS{9}Esse é uma figura para o tempo presente, em que são oferecidos ofertas e sacrifícios que não podem, quanto a consciência, tornar perfeito a quem faz o serviço;
\VS{10}{\ADD{pois consistem}} somente em comidas e bebidas, e vários lavamentos, e ordenanças para o corpo,
\FTNT{‡}{{\FR 9:10 }
corpo
lit. carne – também no v. 13
} que foram impostas até o tempo da correção.
\VS{11}Mas quando veio Cristo, o Sumo Sacerdote dos bens futuros, por meio de um Tabernáculo maior e mais perfeito, não feito por mãos, isto é, não desta criação;
\VS{12}ele entrou de uma vez por todas no Santuário,
\FTNT{§}{{\FR 9:12 }
Há controvérsia se a referência é ao lugar Santo, ou ao Santo dos Santos
} e obteve uma redenção eterna, não pelo sangue de bodes e bezerros, mas sim, pelo seu próprio sangue.
\VS{13}Pois, se o sangue de touros e bodes, e as cinzas de uma novilha espalhadas sobre os imundos, santifica para a purificação do corpo,
\VS{14}quanto mais o sangue do Cristo que, pelo Espírito eterno, ofereceu a si mesmo, imaculado, a Deus, purificará a vossa consciência das obras mortas, para servirdes ao Deus vivo!
\VS{15}E por isso ele é o Mediador de um Novo Testamento,
\FTNT{*}{{\FR 9:15 }
Testamento
a mesma palavra no grego significa testamento ou pacto
} a fim de que, com a ocorrência de uma morte para redenção das transgressões sob o primeiro Testamento, os que foram chamados recebam a promessa da herança eterna.
\VS{16}Pois onde há um testamento, é necessário que ocorra a morte do testador;
\VS{17}porque um testamento se confirma nos mortos, visto que não é válido enquanto o testador vive.
\VS{18}Por isso o primeiro
{\ADD{testamento}} também não foi consagrado sem sangue;
\VS{19}Porque, depois de Moisés haver pronunciado a todo o povo todo mandamento segundo a Lei, ele tomou o sangue de bezerros e de bodes, com água, lã purpúrea, e hissopo, e aspergiu, tanto o próprio livro, como todo o povo,
\VS{20}dizendo: Este é o sangue do pacto que Deus ordenou para vós.
{\RQ{Êxodo 24:8
}}
\VS{21}E semelhantemente aspergiu com o sangue o Tabernáculo, e todos os utensílios do serviço de culto.
\VS{22}Segundo a Lei, quase todas as coisas são purificadas com sangue, e sem derramamento de sangue não há perdão de pecados.
\FTNT{†}{{\FR 9:22 }
perdão de pecados
tradicionalmente: remissão
}
\VS{23}Portanto era necessário que os esboços das coisas que estão nos céus fossem purificados com esses
{\ADD{sacrifícios}} ; mas as próprias coisas celestiais, com sacrifícios melhores que esses.
\VS{24}Pois Cristo não entrou num santuário feito por mãos, mera figura do verdadeiro; mas sim, no próprio Céu, para agora comparecer por nós diante da face de Deus.
\FTNT{‡}{{\FR 9:24 }
Santuário
ou: Santos lugares
}
\VS{25}Também não
{\ADD{entrou}} para oferecer a si mesmo muitas vezes, como o sumo sacerdote entra a cada ano no santuário com sangue alheio;
\VS{26}(de outra maneira lhe seria necessário padecer muitas vezes desde a fundação do mundo), mas agora, no fim dos tempos, ele se manifestou de uma vez por todas para aniquilar o pecado pelo sacrifício de si mesmo.
\VS{27}E, como está ordenado aos seres humanos morrerem uma vez, e depois disso, o juízo,
\VS{28}assim também Cristo, que se ofereceu uma vez para tirar os pecados de muitos, aparecerá pela segunda vez, sem pecado, aos que o esperam, para a salvação.

\par }\Chap{10}{\PP \VerseOne{1}Pois, como a Lei tem uma sombra dos bens futuros, e não a própria imagem das coisas, ela nunca pode, por meio dos mesmos sacrifícios que se oferecem a cada ano, continuamente, tornar perfeitos os que se aproximam.
\FTNT{*}{{\FR 10:1 }
os que se aproximam
i.e., os que se aproximam para oferecer sacrifícios
}
\VS{2}Caso contrário, deixariam de ser oferecidos, pois os adoradores, uma vez purificados, não teriam mais consciência alguma de pecados.
\VS{3}Porém nesses
{\ADD{sacrifícios}} a cada ano
{\ADD{se faz}} uma nova lembrança dos pecados,
\VS{4}porque é impossível que o sangue de touros e de bodes tire pecados.
\VS{5}Por isso que, quando ele entrou no mundo, disse: Sacrifício e oferta não quiseste, mas me preparaste um corpo.
\VS{6}Ofertas de queima e ofertas pelo pecado não te agradaram.
\VS{7}Então eu disse: “Eis-me aqui; (no rolo do livro está escrito de mim) venho para fazer a tua vontade, ó Deus”.
{\RQ{Salmos 40:6-8
}}
\VS{8}Depois de ter dito acima: “Sacrifício, oferta, holocaustos, e ofertas pelo pecado não quiseste, nem te agradaram”, (os quais se oferecem segundo a Lei),
\VS{9}então disse: “Eis-me aqui, venho para fazer a tua vontade, ó Deus”.
{\ADD{Assim}} , ele cancela
\FTNT{†}{{\FR 10:9 }
cancela
lit. tira
} o primeiro
{\ADD{pacto}} ,para estabelecer o segundo.
\VS{10}Pela sua vontade somos santificados por meio da oferta sacrificial do corpo de Jesus Cristo,
{\ADD{feita}} de uma vez por todas.
\VS{11}Todo sacerdote comparece a cada dia para servir e oferecer muitas vezes os mesmos sacrifícios, que nunca podem tirar os pecados;
\VS{12}Mas
{\ADD{Jesus}} ,depois que ofereceu um sacrifício pelos pecados, assentou-se para sempre à direita de Deus;
\VS{13}e espera desde então, até que os seus inimigos sejam postos por escabelo dos seus pés.
\VS{14}Pois com uma só oferta ele aperfeiçoou para sempre os que são santificados.
\VS{15}E também o Espírito Santo nos dá testemunho
{\ADD{acerca disso}} ; pois, depois de haver anteriormente dito:
\VS{16}Este é o pacto que farei com eles depois daqueles dias, diz o Senhor: Porei minhas leis em seus corações, e as escreverei em suas mentes;
{\RQ{Jeremias 31:33
}}
\VS{17}{\ADD{Então ele diz}} : E não mais me lembrarei dos seus pecados e das suas transgressões.
{\RQ{Jeremias 31:34
}}
\VS{18}Ora, onde há perdão
\FTNT{‡}{{\FR 10:18 }
perdão
tradicionalmente: remissão
} dessas coisas, não há mais oferta pelo pecado.
\VS{19}Portanto, irmãos, já que temos a confiança de entrar no Santuário
\FTNT{§}{{\FR 10:19 }
Há controvérsia se a referência é ao lugar Santo, ou ao Santo dos Santos
} pelo sangue de Jesus,
\VS{20}pelo caminho novo e vivo que ele consagrou para nós através do véu, isto é, pela sua carne;
\VS{21}e já que temos um grande Sacerdote sobre a casa de Deus,
\VS{22}aproximemo-nos com um coração sincero em plena certeza de fé, tendo os corações aspergidos e purificados da má consciência, e o corpo lavado com água pura;
\VS{23}mantenhamos firme a esperança que declararmos ter, sem abalo algum, pois aquele que prometeu é fiel;
\FTNT{*}{{\FR 10:23 }
a esperança que declaramos ter
lit. a declaração de fé da [nossa] esperança
}
\VS{24}e sejamos atenciosos uns para com os outros, a fim de incentivar o amor e as boas obras;
\VS{25}não abandonando a nossa reunião, como é o costume de alguns. Ao contrário, encorajemos
\FTNT{†}{{\FR 10:25 }
encorajemos
ou: exortemos
}
{\ADD{uns aos outros}} , e tanto mais quanto vedes aquele dia
\FTNT{‡}{{\FR 10:25 }
aquele dia
lit. o dia
} se aproximando.
\VS{26}Pois se nós, depois de havermos recebido o conhecimento da verdade, persistirmos pecando por vontade própria, já não resta mais sacrifício pelos pecados;
\VS{27}em vez disso, certa terrível expectativa de julgamento, e um fogo de indignação
\FTNT{§}{{\FR 10:27 }
fogo de indignação
lit. zelo de fogo; isto é, a indignação causada pelo zelo
} que consumirá os adversários.
\VS{28}Se alguém que rejeita a Lei de Moisés morre sem misericórdia com base
{\ADD{na palavra de}} duas ou três testemunhas,
\VS{29}de quanto pior castigo vós pensais que será julgado merecedor aquele que pisou o Filho de Deus, menosprezou
\FTNT{*}{{\FR 10:29 }
menosprezou
lit. considerou profano (ou banal
} ) o sangue do Testamento no qual ele foi santificado, e insultou o Espírito da graça?
\VS{30}Pois nós conhecemos aquele que disse: A vingança é minha; eu retribuirei, diz o Senhor. E outra vez: O Senhor julgará o seu povo.
{\RQ{Deuteronômio 32:35,36
}}
\VS{31}Cair nas mãos do Deus vivo é algo terrível.
\VS{32}Lembrai-vos, porém, dos dias passados, em que, depois de serdes iluminados, suportastes grande combate de aflições,
\VS{33}quando, em parte, fostes expostos em público tanto a insultos, como a tribulações, e em parte, fostes companheiros daqueles que assim foram tratados.
\VS{34}Pois também vos compadecestes das minhas prisões, e com alegria aceitastes a espoliação dos vossos bens, pois sabeis em vós mesmos que tendes nos Céus um bem melhor e permanente.
\VS{35}Portanto, não rejeiteis a vossa confiança, que tem uma grande recompensa;
\VS{36}pois precisais de paciência, a fim de que, depois que houverdes feito a vontade de Deus, recebais o que foi prometido.
\FTNT{†}{{\FR 10:36 }
o que foi prometido
lit. a promessa
}
\VS{37}Pois ainda um pouquinho de tempo, e aquele que vem virá, e não tardará.
{\RQ{Isaías 26:20, Habacuque 2:3
}}
\VS{38}Mas o justo viverá pela fé; e se ele retroceder, a minha alma não tem prazer nele.
{\RQ{Habacuque 2:4
}}
\VS{39}Mas nós não somos dos que retrocedem para a perdição, mas sim dos que creem para a conservação da alma.

\par }\Chap{11}{\PP \VerseOne{1}Ora, a fé é a certeza das coisas que se esperam, e a prova das coisas que não se veem.
\VS{2}Pois por ela os antigos obtiveram testemunho.
\FTNT{*}{{\FR 11:2 }
testemunho
ou: aprovação
}
\VS{3}Pela fé entendemos que o universo
\FTNT{†}{{\FR 11:3 }
o universo
lit. as eras – o termo grego literalmente indicava longos períodos de tempo, mas também podia indicar lugares, isto é, o mundo ou mundos
} foi aprontado pela palavra de Deus, de maneira que as coisas que se veem não foram feitas de algo visível.
\VS{4}Pela fé, Abel ofereceu a Deus melhor sacrifício do que Caim; por isso obteve testemunho de que era justo, pois Deus testemunhou de suas ofertas; e mesmo estando morto, ainda fala por meio dela.
\VS{5}Pela fé, Enoque foi transportado para não experimentar
\FTNT{‡}{{\FR 11:5 }
experimentar
lit. ver
} a morte; e não foi achado, porque Deus havia o transportado; pois antes da sua transportação ele obteve testemunho de ter agradado a Deus.
\VS{6}Ora, sem fé é impossível agradar
{\ADD{a Deus}} . Pois é necessário que aquele que se aproxima de Deus creia que ele existe, e que ele é recompensador dos que o buscam.
\VS{7}Pela fé, Noé, divinamente advertido das coisas que ainda se não viam, temeu; e, para salvamento de sua família construiu a arca. Por meio da fé
\FTNT{§}{{\FR 11:7 }
Por meio da fé
lit. Por meio dela
} ele condenou o mundo, e foi feito herdeiro da justiça segundo a fé.
\VS{8}Pela fé, Abraão, quando foi chamado para ir ao lugar que havia de receber por herança, obedeceu, e saiu sem que soubesse aonde ia.
\VS{9}Pela fé, ele habitou na terra de promessa, como em
{\ADD{terra}} que não fosse sua, morando em tendas com Isaque e com Jacó, herdeiros com ele da mesma promessa.
\VS{10}Pois ele esperava a cidade que tem fundamentos firmes, da qual o arquiteto e construtor é Deus.
\VS{11}Pela fé, também, a própria Sara recebeu a capacidade de conceber descendência, e já na idade avançada ela deu à luz, pois considerou fiel aquele que havia prometido.
\FTNT{*}{{\FR 11:11 }
O versículo também pode traduzido se referindo a Abraão: “Pela fé, com Sara, ele recebeu a capacidade de produzir descendência, e já na idade avançada ela deu à luz, pois ele considerou fiel aquele que havia prometido”
}
\VS{12}Assim, de um que estava à beira da morte, nasceu
{\ADD{uma descendência}} numerosa como as estrelas do céu, e incontável como a areia na praia do mar.
\VS{13}{\ADD{Permanecendo}} na fé, todos esses morreram sem receberem as promessas. Mas as viram de longe, creram, e felicitaram-se nelas, declarando que eram estrangeiros e peregrinos na terra.
\VS{14}Pois os que dizem tais coisas mostram claramente que estão buscando a
{\ADD{sua}} própria pátria.
\VS{15}Ora, se estivessem pensando naquela de onde saíram, eles tinham oportunidade de voltar.
\VS{16}Mas agora desejam uma melhor, isto é, a celestial. Por isso Deus não se envergonha deles, de ser chamado o Deus deles, porque já lhes preparou uma cidade.
\VS{17}Pela fé, Abraão, quando foi provado, ofereceu Isaque; aquele que havia recebido as promessas ofereceu o
{\ADD{seu}} único filho,
\VS{18}do qual foi dito: Em Isaque será chamada a tua descendência.
\FTNT{†}{{\FR 11:18 }
lit. semente
}
{\RQ{Gênesis 21:12
}}
\VS{19}{\ADD{Abraão}} considerou que Deus era poderoso até para ressuscitar
{\ADD{Isaque}} dos mortos; de onde também, de forma figurada, ele o recuperou.
\VS{20}Pela fé, Isaque abençoou Jacó e Esaú quanto às coisas futuras.
\VS{21}Pela fé, Jacó, enquanto morria, abençoou a cada um dos filhos de José, e adorou
{\ADD{apoiado}} sobre a ponta do seu cajado.
\VS{22}Pela fé, José, no fim de sua vida,
\FTNT{‡}{{\FR 11:22 }
no fim de sua vida
lit. finalizando [sua vida],isto é, morrendo
} mencionou a saída dos filhos de Israel, e deu ordem acerca dos seus ossos.
\VS{23}Pela fé, Moisés, quando nasceu, foi escondido três meses pelos seus pais, pois viram que o menino era belo, e não se intimidaram com o mandamento do rei.
\VS{24}Pela fé, Moisés, sendo já crescido, recusou ser chamado filho da filha de faraó;
\VS{25}pois preferiu ser maltratado com o povo de Deus a ter por um tempo o prazer do pecado.
\VS{26}Ele considerou que as humilhações por causa de Cristo eram riquezas maiores que os tesouros no Egito; pois sua atenção estava fixada na recompensa.
\VS{27}Pela fé, ele deixou o Egito sem temer a ira do rei, pois perseverou como que vendo aquele que é invisível.
\VS{28}Pela fé, ele celebrou a Páscoa e o derramamento do sangue, para que o destruidor dos primogênitos não os tocasse.
\VS{29}Pela fé, atravessaram o mar Vermelho, como que por
{\ADD{ terra
}} seca, o que, quando os Egípcios tentaram, afogaram-se.
\VS{30}Pela fé, os muros de Jericó caíram, depois de serem rodeados por sete dias.
\VS{31}Pela fé, Raabe, a prostituta, não pereceu com os incrédulos,
\FTNT{§}{{\FR 11:31 }
incrédulos
ou: desobedientes
} pois acolheu em paz os espias.
\VS{32}E que mais direi? Pois me faltará tempo se eu contar a respeito de Gideão, Baraque, e também de Sansão, Jefté, Davi, e também de Samuel e dos profetas.
\VS{33}Pela fé, esses venceram reinos, exercitaram justiça, alcançaram promessas,
\FTNT{*}{{\FR 11:33 }
alcançaram promessas
isto é, alcançaram o que lhes fora prometido
} fecharam bocas de leões,
\VS{34}apagaram o poder do fogo, escaparam do fio de espada, da fraqueza tiraram forças, tonaram-se fortes na batalha, puseram em fuga os exércitos dos estrangeiros;
\VS{35}As mulheres receberam de volta os seus mortos pela ressurreição; e outros foram torturados, não aceitando a soltura, para alcançarem uma melhor ressurreição;
\VS{36}e outros experimentaram provação de insultos e açoites, e até correntes e prisões;
\VS{37}foram apedrejados, serrados, tentados, morreram pela espada; andaram
{\ADD{vestidos}} de peles de ovelhas e de cabras; necessitados, afligidos, maltratados;
\VS{38}(o mundo não era digno deles) andando sem rumo em desertos, montanhas, covas, e nas cavernas da terra.
\VS{39}E todos esses,
{\ADD{mesmo}} tendo testemunho pela fé, não alcançaram a promessa,
\VS{40}pois Deus havia preparado algo melhor para nós, para que, sem nós, eles não fossem aperfeiçoados.
\FTNT{†}{{\FR 11:40 }
sem nós, … aperfeiçoados
isto é: para que eles fossem aperfeiçoados somente conosco
}

\par }\Chap{12}{\PP \VerseOne{1}Portanto nós também, posto que estamos rodeados por uma tão grande nuvem de testemunhas, deixemos toda sobrecarga, e o pecado que facilmente nos envolve, e corramos com perseverança a corrida que nos está proposta,
\VS{2}olhando para Jesus, Autor e aperfeiçoador da fé. Ele, pela alegria que lhe estava proposta, suportou a cruz, desprezando a humilhação, e assentou-se à direita do trono de Deus.
\VS{3}Considerai, pois, aquele que suportou tal hostilidade dos pecadores contra si mesmo, para que não fiqueis exaustos, nem vossas almas se debilitem.
\VS{4}Enquanto lutais contra o pecado, ainda não resististes ao ponto de ter o próprio sangue derramado,
\FTNT{*}{{\FR 12:4 }
resististes ao ponto de ter o próprio sangue derramado
Lit. resististes ao sangue, i.e. resististes a ponto de morrer
}
\VS{5}mas já vos esquecestes do encorajamento que ele fala convosco como a filhos: Meu filho, não desprezes a disciplina do Senhor, nem te canses de ser reprendido por ele;
\VS{6}pois o Senhor disciplina a quem ama, e açoita a todo filho a quem recebe.
{\RQ{3:11-12
}}
\VS{7}Se suportais a disciplina, Deus vos trata como filhos; pois que filho há a quem o pai não discipline?
\VS{8}Mas se estais sem disciplina, da qual todos são feitos participantes, então sois ilegítimos, e não filhos.
\VS{9}Além disso, tivemos os pais de nossa carne como disciplinadores, e nós os respeitávamos. Por acaso não nos sujeitaremos muito mais ao Pai dos espíritos, a fim de vivermos?
\VS{10}Porque eles, por um pouco de tempo,
{\ADD{nos}} disciplinavam como bem lhes parecia. Ele, porém,
{\ADD{disciplina}} para o
{\ADD{nosso}} proveito, a fim de que sejamos participantes da sua santidade.
\VS{11}De fato, no presente, nenhuma disciplina parece ser motivo de alegria, mas sim, de sofrimento; mas depois produz um fruto pacífico de justiça aos que foram exercitados por ela.
\VS{12}Portanto, levantai as mãos cansadas e os joelhos enfraquecidos,
\VS{13}e “fazei caminhos retos para os vossos pés”, para que o que está manco não se desvie, em vez disso, seja sarado.
{\RQ{Provérbios 4:26
}}
\VS{14}Buscai a paz com todos, e a santificação, sem a qual ninguém verá o Senhor.
\VS{15}Vigiai com empenho para que ninguém perca de receber da graça de Deus; para não acontecer que alguma raiz de amargura brote, seja incômoda, e muitos sejam contaminados por ela.
\VS{16}Ninguém seja pecador sexual, ou profano como Esaú, que por uma refeição vendeu o seu direito de primogenitura.
\VS{17}Pois vós sabeis que depois, quando ele quis herdar a benção, foi rejeitado, porque não achou lugar para arrependimento, ainda que com lágrimas o tenha buscado.
\VS{18}Porque não chegastes a um monte palpável, aceso com fogo, e à escuridão, às trevas, e à tempestade;
\VS{19}ao som da trombeta, e à voz das palavras, que os que a ouviam, rogaram que não mais se lhes falasse palavra alguma;
\VS{20}(pois não podiam suportar o que
{\ADD{lhes}} era ordenado: Se até um animal tocar o monte, seja apedrejado ou perfurado com um dardo.
{\RQ{Êxodo 19:12,13
}}
\VS{21}E a visão era tão terrível, que Moisés disse: Estou assombrado e tremendo).
{\RQ{Deuteronômio 9:19
}}
\VS{22}Mas vós chegastes ao monte Sião, à cidade do Deus vivo, a Jerusalém celestial, e aos muitos milhares de anjos,
\VS{23}à universal congregação e igreja dos primogênitos inscritos nos Céus, e a Deus, o Juiz de todos, e aos espíritos dos justos já aperfeiçoados;
\VS{24}e a Jesus, o Mediador de um Novo Testamento, e ao sangue da aspersão, que fala melhores coisas que
{\ADD{o de}} Abel.
\VS{25}Tende o cuidado de não rejeitardes ao que fala. Pois, se não escaparam aqueles que rejeitaram ao que
{\ADD{os}} advertia divinamente na terra, muito menos nós, se nos desviarmos daquele que
{\ADD{nos adverte}} dos Céus.
\VS{26}Sua voz fez a terra tremer naquela ocasião, mas agora ele prometeu, dizendo: Ainda uma vez faço tremer não somente a terra, mas também o Céu.
{\RQ{Ageu 2:6
}}
\VS{27}Esta
{\ADD{expressão}} , “Ainda uma vez”, mostra a remoção das coisas abaladas, como coisas criadas, para que as inabaláveis permaneçam.
\VS{28}Por isso, já que recebemos um Reino inabalável, mantenhamos a graça,
\FTNT{†}{{\FR 12:28 }
graça
ou: gratidão
} e por ela sirvamos a Deus de maneira que o agrade, com devoção e temor,
\VS{29}pois o nosso Deus é um fogo consumidor.
{\RQ{Deuteronômio 4:24
}}

\par }\Chap{13}{\PP \VerseOne{1}Permaneça o amor fraternal.
\VS{2}Não vos esqueçais de mostrar hospitalidade, porque através dela alguns, sem saber, acolheram anjos.
\VS{3}Lembrai-vos dos prisoneiros, como se estivésseis presos com eles; e dos que são maltratados, como se vós mesmos também estivessem sendo em
{\ADD{vossos}} corpos.
\VS{4}O matrimônio seja honrado entre todos, e o leito conjugal sem contaminação; mas Deus julgará os pecadores sexuais e os adúlteros.
\VS{5}A
{\ADD{vossa}} maneira de viver seja sem ganância, contentando-vos com as coisas que tendes. Pois ele disse: Não te deixarei, nem te desampararei.
{\RQ{Deuteronômio 31:6
}}
\VS{6}De maneira que devemos ter a confiança de dizer: O Senhor é meu ajudador; não temerei o que o ser humano poderá fazer a mim.
{\RQ{Salmos 118:6
}}
\FTNT{*}{{\FR 13:6 }
Ou: O Senhor é meu ajudador, não temerei. O que poderá o ser humano fazer a mim?
}
\VS{7}Lembrai-vos de vossos líderes, que vos falaram a palavra de Deus. Observai o resultado da maneira como viveram, e imitai a fé deles.
\VS{8}Jesus Cristo é o mesmo ontem, hoje, e eternamente.
\VS{9}Não vos deixeis levar por doutrinas várias e estranhas. Pois bom ao coração é ser fortificado pela graça, e não por alimentos, que não resultaram em proveito algum aos que com eles se ocuparam.
\VS{10}Temos um altar, do qual os que servem no tabernáculo não têm autoridade para comer;
\VS{11}porque os corpos dos animais, cujo sangue pelo pecado é trazido ao Santuário
\FTNT{†}{{\FR 13:11 }
Há controvérsia se a referência é ao lugar Santo, ou ao Santo dos Santos
} pelo Sumo sacerdote, são queimados fora do acampamento.
\VS{12}Por isso também Jesus, para santificar o povo pelo seu próprio sangue, sofreu fora do portão da cidade.
\VS{13}Assim, saiamos até ele, fora do acampamento, carregando a sua humilhação.
\VS{14}Pois não temos aqui cidade permanente, mas buscamos a futura.
\VS{15}Portanto, por meio dele, ofereçamos continuamente sacrifício de louvor a Deus, isto é, o fruto dos lábios que declaram honra
\FTNT{‡}{{\FR 13:15 }
declaram honra
tradicionalmente: confessam
} ao seu nome.
\VS{16}E não vos esqueçais de fazer o bem e de compartilhar, pois Deus se agrada com tais sacrifícios.
\VS{17}Obedecei aos vossos líderes, e sede submissos
{\ADD{a eles}} ,pois eles vigiam pelas vossas almas, como os que prestarão contas
{\ADD{delas}} ; para que o façam com alegria, e não gemendo, porque isso não vos seria proveitoso.
\VS{18}Orai por nós; porque confiamos que temos boa consciência, e desejamos nos comportar da maneira correta em tudo.
\VS{19}E eu vos rogo ainda mais que façais isto, para que eu vos seja restituído mais depressa.
\VS{20}O Deus da paz, que, pelo sangue do eterno Testamento eterno, voltou a trazer dentre os mortos o grande Pastor das ovelhas, o nosso Senhor Jesus,
\VS{21}vos aperfeiçoe em toda boa obra, para fazerdes a sua vontade, operando em vós o que é agradável diante dele, por meio de Jesus Cristo, a quem seja a glória para todo o sempre. Amém!
\VS{22}Eu vos rogo, porém, irmãos, que suportai esta palavra de exortação; pois eu vos escrevi de maneira breve.
\VS{23}Sabei que o irmão Timóteo já está solto. Se ele vier depressa, com ele eu vos verei.
\VS{24}Saudai a todos os vossos líderes, e a todos os santos. Os da Itália vos saúdam.
\VS{25}A graça seja com todos vós. Amém.{\ADD{Escrita da Itália para os Hebreus, e enviada por Timóteo}}
\par }